% Section: Limitations
\section{Limitations}
\label{sec:limitations}

\paragraph{Evaluation scope.}
Our strongest empirical gains appear on the routing benchmark, where all four strategies are
executed offline per query and oracle correctness and cost labels are available for training
and regret evaluation.
Benchmark transfer is a secondary evaluation with mixed exact-match outcomes (e.g., pooled
48.1\% vs.\ 50.3\% for a heuristic baseline); the proposed router does not uniformly lead on
every executed family despite competitive performance.
\textbf{Consequently}, routing-benchmark regret improvements should not be read as uniform
accuracy gains on every public workload without per-family transfer validation.
The framework is most reliable where training and evaluation share oracle multi-strategy
supervision.

\paragraph{Fixed strategy pool and shared backbone.}
We route among four procedurally distinct strategies implemented on a shared model backbone.
We do not study factorial combinations of backbones and agent designs, tool inventories, or
prompt variants.
Expanding $\mathcal{A}$ may change both the feature space and the supervision distribution
required for cost-soft targets.

\paragraph{Planner cost and quality.}
Procedure-graph complexity depends on an LLM planner whose outputs can be noisy or miscalibrated
for out-of-domain queries.
When QCE runs at inference, total regret includes planner and feature-extraction cost;
Selective QCE (Appendix~\ref{app:selective}) partially mitigates this overhead but is not the
primary system reported in the main tables.

\paragraph{Language and benchmark coverage.}
Benchmark workloads are English-centric and drawn from reasoning, knowledge, and agentic QA
settings.
Generalization to multilingual inputs, long-context documents, or safety-critical domains is
untested.

\paragraph{Literature baselines.}
RouteLLM, LLM-Blender, and GraphRouter inform our objective, supervision, and feature design but
are not drop-in reimplemented in the four-strategy agent setting because of task mismatch
(model-tier routing vs.\ tool-using agent pipelines).
Direct head-to-head numbers against those systems remain an open comparison when unified
benchmarks become available.
